% Bagian: sets-functions-relations
% Bab: functions
% Subbagian: kinds

\documentclass[../../../include/open-logic-section]{subfiles}

\begin{document}

\olfileid[id]{sfr}{fun}{kin}
\olsection{Jenis-Jenis Fungsi}

\begin{explain}
Akan berguna jika kita memperkenalkan semacam klasifikasi bagi beberapa jenis
fungsi yang paling sering kita jumpai.

Sebagai permulaan, kita dapat meninjau fungsi yang mempunyai sifat bahwa
setiap anggota kodomain merupakan nilai fungsi tersebut. Fungsi semacam itu
disebut !!{surjective}, dan dapat digambarkan seperti dalam
\olref{fig:surjective}.

\begin{figure}
  \olasset{assets/diagrams/surjective.tikz}
  \caption{Fungsi !!^a{surjective} mempunyai setiap !!{element} kodomain
    sebagai suatu nilai.}
  \ollabel{fig:surjective}
\end{figure}
\end{explain}

\begin{defn}[Fungsi !!^{surjective}]
Suatu fungsi $f \colon A \rightarrow B$ bersifat \emph{!!{surjective}} jika
dan hanya jika $B$ juga merupakan daerah hasil dari~$f$, yaitu untuk setiap
$y \in B$ terdapat sekurang-kurangnya satu $x \in A$ sehingga~$f(x) = y$,
atau dalam simbol:
\[
  (\forall y \in B)(\exists x \in A)f(x) = y.
\]
Fungsi semacam itu kita sebut !!a{surjection} dari $A$ ke $B$.
\end{defn}

\begin{explain}
Jika Anda hendak menunjukkan bahwa $f$ merupakan !!a{surjection}, Anda harus
menunjukkan bahwa setiap objek dalam kodomain $f$ adalah nilai $f(x)$ untuk
suatu masukan $x$.

Perhatikan bahwa setiap fungsi \emph{menginduksi} !!a{surjection}. Jika
diberikan suatu fungsi $f \colon A \to B$, definisikan $f' \colon A \to
\ran{f}$ dengan $f'(x) = f(x)$. Karena $\ran{f}$ \emph{didefinisikan} sebagai
$\Setabs{f(x) \in B}{x \in A}$, fungsi $f'$ ini pasti merupakan
!!a{surjection}.
\end{explain}

\begin{explain}
Setiap fungsi memetakan tiap masukan yang mungkin ke satu-satunya keluaran.
Namun, ada pula fungsi yang tidak pernah memetakan masukan yang berbeda ke
keluaran yang sama. Fungsi semacam itu disebut !!{injective}, dan dapat
digambarkan seperti dalam \olref{fig:injective}.
\begin{figure}
  \olasset{assets/diagrams/injective.tikz}
  \caption{Fungsi !!^a{injective} tidak pernah memetakan dua argumen yang
    berbeda ke nilai yang sama.}
  \ollabel{fig:injective}
\end{figure}
\end{explain}

\begin{defn}[Fungsi !!^{injective}]
Suatu fungsi $f \colon A \rightarrow B$ bersifat \emph{!!{injective}} jika
dan hanya jika untuk setiap $y \in B$ terdapat paling banyak satu $x \in A$
sehingga~$f(x) = y$. Fungsi semacam itu kita sebut !!a{injection} dari $A$
ke~$B$.
\end{defn}

\begin{explain}
Jika Anda hendak menunjukkan bahwa $f$ merupakan !!a{injection}, Anda harus
menunjukkan bahwa untuk sebarang !!{element}s $x$ dan $y$ dari domain $f$,
jika $f(x)=f(y)$, maka $x=y$.
\end{explain}

\begin{ex}
Fungsi konstan $f\colon \Nat \to \Nat$ yang diberikan oleh $f(x) = 1$ tidak
bersifat !!{injective} maupun !!{surjective}.

Fungsi identitas $f\colon \Nat \to \Nat$ yang diberikan oleh $f(x) = x$
bersifat !!{injective} sekaligus !!{surjective}.

Fungsi penerus $f \colon \Nat \to \Nat$ yang diberikan oleh $f(x) = x+1$
bersifat !!{injective}, tetapi tidak !!{surjective}.

Fungsi $f \colon \Nat \to \Nat$ yang didefinisikan oleh:
\[
  f(x) =
  \begin{cases}
    \frac{x}{2} & \text{jika $x$ genap} \\
    \frac{x+1}{2} & \text{jika $x$ ganjil.}
  \end{cases}
\]
bersifat !!{surjective}, tetapi tidak !!{injective}.
\end{ex}

\begin{explain}
Cukup sering kita hendak meninjau fungsi yang sekaligus !!{injective} dan
!!{surjective}. Fungsi semacam itu kita sebut !!{bijective}. Bentuknya seperti
fungsi yang digambarkan dalam \olref{fig:bijective}. !!^{bijection}s
kadang-kadang juga disebut \emph{korespondensi satu-satu}, karena memasangkan secara
unik elemen-elemen kodomain dengan elemen-elemen domain.
\begin{figure}
  \olasset{assets/diagrams/bijective.tikz}
  \caption{Fungsi !!^a{bijective} memasangkan secara unik elemen-elemen
    kodomain dengan elemen-elemen domain.}
  \ollabel{fig:bijective}
\end{figure}
\end{explain}

\begin{defn}[!!^{bijection}]
Suatu fungsi $f \colon A \to B$ bersifat \emph{!!{bijective}} jika dan hanya
jika fungsi itu !!{surjective} sekaligus !!{injective}. Fungsi semacam itu kita
sebut !!a{bijection} dari $A$ ke~$B$ (atau antara $A$ dan~$B$).
\end{defn}

\end{document}
